\begin{frame}[noframenumbering,plain]
    \setcounter{framenumber}{1}
    \maketitle
\end{frame}

\section{Постановка задачи}

\begin{frame}
    \frametitle{Проблема и ограничения}
    \begin{itemize}
        \item Почтовый архив содержит рабочую переписку, уведомления, документы
              и~разовые сообщения на~русском и~английском языках.
        \item Ручная сортировка требует постоянных затрат, а~правила
              по~отправителю или~теме быстро устаревают.
        \item Готовые клиенты обычно выделяют общие категории
              <<важное~/ неважное>>, но~не~формируют персональные классы
              пользователя.
        \item Текст писем содержит персональные и~рабочие данные, поэтому
              обработка должна выполняться локально.
    \end{itemize}
\end{frame}

\begin{frame}
    \frametitle{Актуальность работы}
    \begin{itemize}
        \item Объём личных, учебных и~рабочих текстовых документов растёт,
              поэтому ручное упорядочение требует повторяющихся действий.
        \item Поиск по~названию, отправителю или~отдельным словам плохо
              поддерживает личные категории пользователя.
        \item Тексты документов и~писем могут содержать персональные данные
              и~внутренние материалы.
        \item Поэтому система классификации должна работать локально:
              подготовка текста, построение признаков и~назначение категории
              выполняются на~стороне пользователя.
    \end{itemize}
\end{frame}

\begin{frame}
    \frametitle{Цель работы}
    \textbf{Цель:} разработать систему классификации текстовых документов,
    которая работает локально, назначает категории пользователя
    и~позволяет оценить результат классификации.

    \medskip
    Для проверки цели в~работе использованы:
    \begin{itemize}
        \item локальная обработка почтового архива;
        \item ручное подтверждение пользовательских категорий;
        \item метрики классификации, покрытие и~доля отказов.
    \end{itemize}
\end{frame}

\begin{frame}[t]
    \frametitle{Задачи работы}
    \small
    \setlength{\itemsep}{0.25em}
    \begin{enumerate}
        \item Изучить способы автоматической категоризации текстов.
        \item Определить порядок формирования пользовательских категорий
              и~отказа от~низкоуверенных решений.
        \item Подготовить набор данных электронной почты для~построения эмбеддингов.
        \item Разработать схему локальной обработки: подготовку текста, поиск
              похожих документов и~назначение категории.
        \item Реализовать рабочую версию системы и~проверить её.
    \end{enumerate}
\end{frame}

\begin{frame}
    \frametitle{Краткое описание проекта}
    \begin{itemize}
        \item Система получает письма из~почтового ящика, сохраняет исходные
              сообщения и~строит нормализованный текст.
        \item Готовая модель эмбеддингов формирует векторные представления
              писем; похожие письма группируются алгоритмом кластеризации.
        \item Представители кластеров проходят ручную проверку через интерфейс
              разметки, после чего согласованные метки переносятся на~кластеры.
        \item По размеченным письмам строится прототипный классификатор
              с~порогом уверенности и~служебной категорией \texttt{other}.
        \item Проведена экспериментальная проверка на~почтовом архиве:
              рассчитаны метрики классификации, покрытие и~доля отказов.
    \end{itemize}
\end{frame}
